%
\section{\crowsFoot{G}{O1}{H}{MM}}%
%
\begin{frame}[t]%
\frametitle{\crowsFoot{G}{O1}{H}{MM}}%
\begin{itemize}%
\item Es gibt zwei Entitätstypen~G und~H.%
\item<2-> Jede Entität vom Typ~G muss mit mindestens einer aber vielleicht vielen Entitäten vom Typ~H verbunden sein.%
\item<3-> Jede Entität vom Typ~H ist mit keiner oder einer Entität vom Typ~G verbunden.
\item<4-> Beispiel:\uncover<5->{%
\begin{itemize}%
%
\item In einem Fußballklub trainiert jeder Trainer mehrere Mitglieder.
\uncover<6->{ %
Jedes Mitglied kann von einem Trainer trainiert werden oder auch gar nicht trainiert werden.}%
\end{itemize}%
}%
\end{itemize}%
%
\locateGraphic{4-}{width=0.9\paperwidth}{graphics/GH}{0.05}{0.6}%
\end{frame}%
%
\begin{frame}%
\frametitle{\crowsFoot{G}{O1}{H: }{MM}Lösung}%
\parbox{0.435\linewidth}{\small%
\begin{itemize}%
%
\only<-5>{%
\item Wir erstellen wieder zwei Tabellen und nennen sie diesmal \sqlil{g} und~\sqlil{h}.%
}%
%
\only<-6>{%
\item<2-> Sie haben die Primärschlüssel~\sqlil{gid} und~\sqlil{hid}, sowie die Beispielattribute~\sqlil{x} und~\sqlil{y}.%
}%
%
\only<-7>{%
\item<3-> Wir können die Situation so ähnlich anpacken wie die \crowsFoot{E}{O1}{F}{OM}-Situation.%
}%
%
\only<-8>{%
\item<4-> Allerdings ist es schwerer, dieses Beziehungsschema zu erzwingen.%
}%
%
\only<-9>{%
\item<5-> Wir fangen damit an, Tabelle~\sqlil{g} vorzubereiten.%
}%
%
\only<-10>{%
\item<6-> Jede Entität vom Typ~G muss mit mindestens einer Entität vom Typ~H in Beziehung stehen.%
}%
%
\only<-11>{%
\item<7-> Wir lösen das Problem in zwei Schritten.%
}%
%
\only<-12>{%
\item<8-> Erst erzwingen wir, dass jede Zeile in Tabelle~\sqlil{g} mit einer Zeile in Tabelle~\sqlil{h} verbunden ist.%
}%
%
\only<-12>{%
\item<9-> Später erlauben wir, dass sie mit mehr als einer Zeile verbunden ist.%
}%
%
\only<-13>{%
\item<10-> Wir fangen also mit dem Hinzufpgen eines Attributs~\sqlil{fkhid} an, welches \sqlil{NOT NULL} sein muss.%
}%
%
\only<-14>{%
\item<11-> Diese Spalte sollte immer auf die Zeile in Tabelle~\sqlil{h} mit dem selben Wert in~\sqlil{hid} zeigen.%
}%
%
\only<-15>{%
\item<12-> Später setzen wir eine passende \sqlil{REFERENCE}-Einschränkung über~\sqlil{ALTER TABLE}~(noch können wir as nicht, es gibt die Tabelle~\sqlil{h} ja noch nicht).%
}%
%
\only<-15>{%
\item<13-> Stellen Sie sich erstmal vor, wir hätten das schon gemacht und dass das Attribut immer eine Zeile in Tabelle~\sqlil{h} referenziert.%
}%
%
\only<-17>{%
\item<14-> Über dieses Attribut~\sqlil{fkhid} soll immer die \inQuotes{erste} H-Entität referenzieren, zu der die Zeile in Tabelle~\sqlil{g} in Beziehung steht.%
}%
%
\only<-19>{%
\item<15-> Wir machen diese Spalte \sqlil{UNIQUE}, denn keine Zeile in Tabelle~\sqlil{h} kann mit mehr als einer Zeile in Tabelle~\sqlil{g} verbunden sein.%
}%
%
\only<-20>{%
\item<16-> Deshalb kann kein Primärschlüsselwert~\sqlil{hid} mehr als einmal in der Spalte~\sqlil{fkhid} auftauchen.%
}%
%
\only<-21>{%
\item<17-> Nun bearbeiten wir die Tabelle für den Entitätstyp~H.%
}%
%
\only<-21>{%
\item<18-> Weil jede Entität vom Typ~H mit keiner oder einer Entität vom Typ~G in Verbindung steht, fügen wir die Spalte~\sqlil{fkgid} zur Tabelle~\sqlil{h} hinzu.%
}%
%
\only<-22>{%
\item<19-> Die Werte in dieser Spalte können \sqlil{NULL} sein.%
}%
%
\only<-23>{%
\item<20-> In diesem Fall ist die entsprechende Zeile in~\sqlil{h} mit keiner Entität vom Typ~G verbunden.%
}%
%
\only<-23>{%
\item<21-> Wenn der Wert nicht \sqlil{NULL} ist, dann muss es ein passender Fremdschlüsselverweis auf eine Zeile in Tabelle~\sqlil{g} sein.%
}%
%
\only<-24>{%
\item<22-> Die Werte in dieser Spalte brauchen nicht \sqlil{UNIQUE} zu sein, weil mehrere Entitäten vom Typ~H mit der selben Entität vom Typ~G in Verbindung stehen können.%
}%
%
\only<-25>{%
\item<23-> Jetzt haben wir sichergestellt, dass jede Entität vom Typ~H mit keiner oder genau einer Entität vom Typ~G verbunden ist.%
}%
%
\only<-27>{%
\item<24-> Wir könnten jetzt eine Einschränkung~\sqlil{g_fkhid_fk} als~\sqlil{FOREIGN KEY (fkhid)} \sqlil{REFERENCES h (hid)} zur Tabelle~\sqlil{g} mit~\sqlil{ALTER TABLE g} \sqlil{ADD CONSTRAINT...} hinzufügen.%
}%
%
\only<-28>{%
\item<25-> Das würde erzwingen, dass jede Entität vom Typ~G mit mindestens einer Entität vom Typ~H verbunden sein muss.%
}%
%
\only<-29>{%
\item<26-> Damit stellen wir aber \emph{nicht} sicher, dass wenn eine Entität vom Typ~G mit einer Entität vom Typ~H verunden ist~(als ihre \inQuotes{erste} H-Entität) {\dots} dass diese H-Entität auch mit der G-Entität verbunden ist.%
}%
%
\only<-31>{%
\item<27-> Wir können das auf eine etwas komische Art sicherstellen.%
}%
%
\only<-31>{%
\item<28-> Wir erstellen eine \sqlil{REFERENCES}-Einschränkung nicht auf \emph{eine einzelne} Spalte~\sqlil{hid}\dots%
}%
%
\only<-32>{%
\item<29-> {\dots}sondern wir erzwingen, dass das Paar~\sqlil{(fkhid, gid)} in Tabelle~\sqlil{g} genauso als Paar~\sqlil{(hid, fkgid)} in Tabelle~\sqlil{h} auftauchen muss!%
}%
%
\only<-33>{%
\item<30-> Wir wissen, dass jeder Wert von~\sqlil{gid} nur einmal in Tabelle~\sqlil{g} existiert.%
}%
%
\only<-34>{%
\item<31-> Wir wissen auch, dass jeder Wert von~\sqlil{hid} nur einmal in Tabelle~\sqlil{h} existiert.%
}%
%
\only<-35>{%
\item<32-> Der erste Teil des Spaltenpaares in der Einschränkung -- \sqlil{fkhid} oder von der anderen Seite kommend,~\sqlil{hid} --  selektiert also immer eine eindeutige Zeile in Tabelle~\sqlil{h}.%
}%
%
\only<-36>{%
\item<33-> Es kann niemals eine andere Zeile mit dem selben \sqlil{hid}-Wert geben, denn das ist der Primärschlüssel der Tabelle~\sqlil{h}.%
}%
%
\only<-36>{%
\item<34-> Das zweite Element des Paares -- also \sqlil{gid} oder von der anderen Seite kommend,~\sqlil{fkgid}) -- erzwingt darum, dass diese einzelne Zeile in Tabelle~\sqlil{h} den passenden Wert~\sqlil{gid} in~\sqlil{fkgid} gespeichert hat.%
}%
%
\only<-37>{%
\item<35-> Weil Fremdschlüssel-\sqlil{REFERENCES}-Einschränkungen nur Spalten referenzieren können, die \sqlil{UNIQUE} sind, müssen wir die Einschränkung~\sqlil{UNIQUE (hid, fkgid)} zur Tabelle~\sqlil{h} hinzufügen.%
}%
%
\only<-38>{%
\item<36-> Überlegen wir uns das noch einmal.%
}%
%
\only<-39>{%
\item<37-> Wir haben die Einschränkung \sqlil{g_fkhid_gid_fk} als \sqlil{FOREIGN KEY (fkhid, gid)} \sqlil{REFERENCES h (hid, fkgid)}.%
}%
%
\only<-41>{%
\item<38-> Auf der Seite der Tabelle~\sqlil{g} schaut diese Einschränkung auf eine Zeile und nimmt den Wert des Fremdschlüssels~\sqlil{fkhid} zur Tabelle~\sqlil{h} zusammen mit dem eigenen Primärschlüssel~\sqlil{gid} als ein Tuple~\sqlil{(fkhid, gid)}.%
}%
%
\only<-42>{%
\item<39-> Auf Seiten der Tabelle~\sqlil{h} muss es ein passendes Tupel mit dem Primärschlüssel~\sqlil{hid} und dem dazugehörigen Wert des Fremdschlüssels~\sqlil{fkgid} geben.%
}%
%
\only<-44>{%
\item<40-> Es wird also erzwungen, dass die Paare \sqlil{(fkhi, gid) == (hid, fkgid)} immer passen.%
}%
%
\only<-45>{%
\item<41-> Natürlich sind die Primärschlüsselwerte beider Tabellen immer einzigartig.%
}%
%
\only<-46>{%
\item<42-> Sagen wir, Tabelle~\sqlil{g} hat Primärschlüsselwert~\sqlil{gid=u} und Fremdschlüssel~\sqlil{fkhid=v}.%
}%
%
\item<43-> Dann muss die Zeile in Tabelle~\sqlil{h} mit Primärschlüsselwert~\sqlil{hid=v} den Fremdschlüssel~\sqlil{fkgid=u} haben.%
%
\item<44-> Natürlich kann es auch andere Zeilen in Tabelle~\sqlil{h} geben, die auch den Fremdschlüssel~\sqlil{fkgid=u} haben.%
%
\item<45-> Das ist OK, denn mehrere Entitäten vom Typ~H können die selbe Entität vom Typ~G referenzieren.%
%
\item<46-> Jetzt haben wir also diese Beziehung exakt implementiert.%
%
\item<47-> Schauen wir uns mal an, was wir gemacht haben.%
%
\end{itemize}%
}%
\gitLoadAndExecSQL{GH_tables}{}{conceptualToRelational}{GH_tables.sql}{relationships}{}{}%
\listingSQLandOutput{}{GH_tables}{}{0.47}{0.1}{0.529}{0.87}
\end{frame}%
%
%
\begin{frame}[t]%
\frametitle{\crowsFoot{G}{O1}{H: }{MM}Testen der Einschränkungen}%
\parbox{0.435\linewidth}{\small%
\begin{itemize}%
%
\only<-5>{%
\item Probieren wir mal aus, was wir da gemacht haben.%
}%
\only<-6>{%
\item<2-> Können wir eine Zeile in Tabelle~\sqlil{g} einfügen, die nicht mit einer Zeile in Tabelle~\sqlil{h} in Beziehung steht?%
}%
\only<-7>{%
\item<3-> Nein, können wir nicht.%
}%
\only<-8>{%
\item<4-> Damit wir eine Zeile in Tabelle~\sqlil{g} einfügen können, müssen wir den Fremdschlüssel~\sqlil{fkhid} angeben und die entsprechende Zeile in Tabelle~\sqlil{h} muss es geben.%
}%
%
\item<5-> Können wir Zeilen in Tabelle~\sqlil{h} einfügen, die nicht mit Zeilen in Tabelle~\sqlil{g} in Beziehung stehen?%
%
\item<6-> Ja, das würde gehen. Und das ist OK.%
%
\item<7-> Entitäten vom Typ~H sind mit entweder einer oder keiner Entität vom Tyb~G verbunden.%
%
\item<8-> Aber können wir eine Zeile in Tabelle~\sqlil{g} einfügen, die eine Zeile in Tabelle~\sqlil{h} referenziert, die \emph{nicht} mit einer Entität vom Typ~G verbunden ist?%
%
\item<9-> Nein, denn die Einschränkung~\sqlil{g_fkhid_gid_fk} verlangt, dass die Zeile in Tabelle~\sqlil{h} zurück auf die Zeile in Tabelle~\sqlil{g} verweist.%
%
\end{itemize}%
}%
%
\gitLoadAndExecSQL{GH_insert_error_1}{}{conceptualToRelational}{GH_insert_error_1.sql}{relationships}{}{}%
\gitLoadAndExecSQL{GH_insert_error_2}{}{conceptualToRelational}{GH_insert_error_2.sql}{relationships}{}{}%
\listingSQLandOutput{2-4}{GH_insert_error_1}{}{0.47}{0.1}{0.529}{0.87}
\listingSQLandOutput{8-9}{GH_insert_error_2}{}{0.47}{0.1}{0.529}{0.87}
\end{frame}%
%
\begin{frame}%
\frametitle{\crowsFoot{G}{O1}{H: }{MM}Wie können wir überhaupt Daten einfügen?}%
\begin{itemize}%
\item Aber wie können wir überhaupt Zeilen in Tabelle~\sqlil{g} einfügen?%
%
\item<2-> Wie können keine Zeile einfügen, die keine existierende Zeile in Tabelle~\sqlil{h} referenziert.%
%
\item<3-> Wir können auch keine Zeile einfügen, die gar eine existierende Zeile in Tabelle~\sqlil{h} referenziert, die noch mit keiner Zeile in Tabelle~\sqlil{g} verbunden ist.%
%
\item<4-> Und wenn wir eine Zeile in Tabelle~\sqlil{g} einfügen wollen würden, die eine Zeile in Tabelle~\sqlil{h} referenziert welche bereits eine andere Zeile in Tabelle~\sqlil{q} referenziert {\dots} dann müsste diese andere Zeile in Tabelle~\sqlil{g} erstmal existieren.%
%
\item<5-> Was sie nicht tut.%
%
\item<6-> Nun haben wir also unsere beiden Tabellen erzeugt und ihre referenzielle Integrität mit starken Einschränkungen geschützt{\dots}%
%
\item<7-> {\dots}und die \emph{müssen} ja auch so sein, denn wir wollen ja gerade \crowsFoot{G}{O1}{H}{MM} implementieren
%
\item<8->{\dots}aber wie bekommen wir überhaupt Daten in Tabelle~\sqlil{g}?%
%
\item<9-> Wir müssen dieses komische Henne-Ei-Problem lösen.%
\end{itemize}%
\end{frame}%
%
%
\begin{frame}[t]%
\frametitle{\crowsFoot{G}{O1}{H: }{MM}Befüllen mit Daten}%
\parbox{0.435\linewidth}{\small%
\begin{itemize}%
\only<-4>{%
\item Wir fangen damit an, erstmal Zeilen in die Tabelle~\sqlil{h} zu schreiben.%
}%
%
\only<-6>{%
\item<2-> Da die Entitäten vom Typ~H nicht unbedingt mit denen vom Typ~G in Beziehung stehen müssen, brauchen wir uns um die Einschränkungen keine Gedanken zu machen.%
}%
%
\only<-7>{%
\item<3-> Wenn wir aber eine Entität vom Typ~G in die Tabelle~\sqlil{g} einfügen, dann müssen wir \alert{gleichzeitig} eine Bezihung zu einer H-Entität in der Tabelle~\sqlil{h} erstellen.%
}%
%
\only<-8>{%
\item<4-> Das klingt unmöglich, geht aber problemlos\only<-4>{.}\uncover<5->{, weil:}%
}%
%
\only<6->{%
\only<-9>{%
\item<6-> In \postgresql\ und wahrscheinlich vielen anderen \glsShort{dbms}en werden Einschränkungen zur referentiellen Integrität am Ende einer Transaktion überprüft.%
}%
%
\only<-10>{%
\item<7-> Eine Transaktion kann mehrere Kommandos gruppieren, die dann entweder alle zusammen erfolgreich sind oder alle zusammen fehlschlagen~(wobei dann die \glsShort{db} nicht verändert wird).%
}%
%
\only<-11>{%
\item<8-> Ein einzelnes Kommando in \postgresql\ ist auch eine Transaktion\cite{PGDG:T}.%
}%
%
\only<-12>{%
\item<9-> Das heist, dass Einschränkungen zur referentiellen Integriät erst \emph{am Ende} des Kommandos überprüft werden.%
}%
%
\only<-13>{%
\item<10-> Die Änderungen werden dann zur Datenbank committed, wenn alle Einschränkungen \inQuotes{passen} und ansonsten zurückgerollt.%
}%
%
\only<-14>{%
\item<11-> Um nun eine Zeile in Tabelle~\sqlil{g} einzufügen, dann müssen wir auch einen Datensatz in Tabelle~\sqlil{h} verändern, der aktuell noch auf keine Zeile in Tabelle~\sqlil{g} verweist.%
}%
%
\only<-15>{%
\item<12-> Wir müssen diese Zeile in Tabelle~\sqlil{h} dann mit der neuen Zeile in Tabelle~\sqlil{g} verbinden.%
}%
%
\only<-16>{%
\item<13-> Wenn wir die Einfügung und die Veränderung in ein einziges \glsShort{SQL}-Kommando packen können, dann wird es einfach.%
}%
%
\only<-17>{%
\item<14-> Sonst müssen wir eben lernen, wie man explizit Transaktionen verwendet. Das geht auch.%
}%
%
\only<-18>{%
\item<15-> Um eine den \sqlil{fkgid}-Wert einer Zeile in Tabelle~\sqlil{h} zu verändern, müssen wir den Wert des Primärschlüssels~\sqlil{gid} der neuen Zeile in Tabelle~\sqlil{g} kennen.%
}%
%
\only<-19>{%
\item<16-> Wir haben den Primärschlüssel~\sqlil{gid} der Tabelle~\sqlil{g} als \sqlil{INT GENERATED} \sqlil{BY DEFAULT} \sqlil{AS IDENTITY} \sqlil{PRIMARY KEY} deklariert.%
}%
%
\only<-20>{%
\item<17-> Das bedeutet, dass der Wert dieses Schlüssels erst erzeugt wird, wenn die Zeile generiert wird.%
}%
%
\only<-21>{%
\item<18-> Das heist wiederum, dass wir den Wert nicht kennen können, bevor wir die Zeile in Tabelle~\sqlil{g} einfügen.%
}%
%
\only<-21>{%
\item<19-> Zum Glück ist das ein sehr normales Problem:~\inQuotes{Was, wenn wir Daten in eine Tabelle mit einem automatisch generierten Primärschlüssel einfügen und wir dann den Schlüssel wissen müssen?}%
}%
%
\only<-22>{%
\item<20-> In \postgresql, das geht mit dem \sqlilIdx{RETURNING}-Schlüsselwort\cite{SE:DA:2020WDSSOAOTNPSDIOR,PGDG:PD:RDFMR}.%
}%
%
\only<-23>{%
\item<21-> Das geht nicht bei allen \glsShort{dbms}en, aber bei \mariadb\cite{M:MSD:IR} und \sqlite\cite{HWACIS:R} geht es auch.%
}%
%
\only<-24>{%
\item<22-> Mit \sqlil{INSERT INTO g} \sqlil{(fkhid, x)} \sqlil{VALUES (1, '123')} \sqlil{RETURNING gid, fkhid} würden wir eine Zeile in Tabelle~\sqlil{g} einfügen, wo der Wert des Attributs~\sqlil{x} gleich~\sqlil{'123'} ist und der Wert des Attributs~\sqlil{fkhid} gleich~\sqlil{1}.%
}%
%
\only<-26>{%
\item<23-> Das Kommando würde den Wert des automatisch erzeugten Primärschlüssels~\sqlil{gid} zurückliefern und auch den Wert von \sqlil{fkhid}, welcher in diesem Fall \sqlil{1} wären.%
}%
%
\only<-28>{%
\item<24-> Das Kommando würde natürlich fehlschlagen, denn es verletzt die Einschränkung \sqlil{g_fkhid_gid_fk}.%
}%
%
\only<-29>{%
\item<25-> Aber wir sind schon mal einen Schritt weiter.%
}%
%
\only<-30>{%
\item<26-> Eine Idee wäre es, zu versuchen eine Anfrage wie folgt zu bauen: \sqlil{UPDATE h} \sqlil{SET gid =} \sqlil{(INSERT INTO g} \sqlil{(fkhid, x)} \sqlil{VALUES (1, '123')} \sqlil{RETURNING gid, fkhid)} \sqlil{WHERE h.hid = fkhid;}.%
}%
%
\only<-31>{%
\item<27-> Das funktioniert aber nicht.%
}%
%
\only<-31>{%
\item<28-> Was aber geht ist ein sogenannter \glsFull{CTE}.%
}%
%
\only<-32>{%
\item<29-> Ein \glsShort{CTE} erlaubt es uns, einem Unterausdruck einer Anfrage einen Namen zuzuweisen.%
}%
%
\only<-32>{%
\item<30-> Dann funktioniert der Unterausdruck in etwa wie eine temporäre Tabelle.%
}%
%
\only<-33>{%
\item<31-> Der Unterausdruck wird genau einmal ausgeführt und kann wie eine read-only Tabelle in den anderen Teilen der Anfrage verwendet werden.%
}%
%
\only<-35>{%
\item<32-> Mit \sqlil{WITH cats AS} \sqlil{(SELECT age, name} \sqlil{FROM animals} \sqlil{WHERE type='cat')} würden wir das Ergebnis der Anfrage \sqlil{SELECT age, name} \sqlil{FROM animals} \sqlil{WHERE type='cat'} dem \glsShort{CTE}~\sqlil{cats} zu weisen.%
}%
%
\only<-37>{%
\item<33-> Wir könnten dann \sqlil{cats} verwenden, so als wäre es eine read-only Tabelle, und \DEzB\ mit \sqlil{SELECT name FROM cats;} die Katzennamen abfragen.%
}%
%
\only<-38>{%
\item<34-> Es ist ein wenig so wie ein \sqlilIdx{VIEW}, aber es ist Teil eines einzelnen \glsShort{SQL}-Kommandos.%
}%
%
\only<-39>{%
\item<35-> Nun setzen wir alles zusammen.%
}%
%
\only<-39>{%
\item<36-> Nehmen wir an, wir haben bereits eine Zeile in Tabelle~\sqlil{h} mit Primärschlüssel~1 eingefügt.%
}%
%
\only<-40>{%
\item<37-> Wir können das Jederzeit, denn diese Zeilen müssen nicht mit den Zeilen in Tabelle~\sqlil{g} verbunden sein.%
}%
%
\only<-41>{%
\item<38-> Dann tun wir das Kommando zum Einfügen einer Zeile in Tabelle~\sqlil{g} in einen \glsShort{CTE}~\sqlil{g_new} in dem wir schreiben \sqlil{WITH g_new AS} \sqlil{(INSERT INTO g} \sqlil{(fkhid, x) VALUES} \sqlil{(1, '123')} \sqlil{RETURNING id, fkhid)}\sqlIdx{WITH}.%
}%
%
\only<-43>{%
\item<39-> Dieser \glsShort{CTE} fügt eine Zeile in Tabelle~\sqlil{g} ein, die die Zeile in Tabelle~\sqlil{h} mit Primärschlüssel~1 referenziert.%
}%
%
\only<-43>{%
\item<40-> Er liefert auch den Primärschlüssel der neuen Zeile in Tabelle~\sqlil{g} als~\sqlil{gid} zurück, ebenso den Fremdschlüssel \sqlil{fkhid}~(mit Wert~1).%
}%
%
\only<-45>{%
\item<41-> Dann benutzen wir diesen \glsShort{CTE} um die Zeile mit Primärschlüssel \sqlil{hid = fkhid} in Tabelle~\sqlil{h} upzudaten.%
}%
%
\only<-46>{%
\item<42-> Wir machen \sqlil{UPDATE h SET} \sqlil{fkgid = g_new.gid} \sqlil{FROM g_new} \sqlil{WHERE h.gid = g_new.fkhid;}.%
}%
%
\only<-47>{%
\item<43-> Dadurch zeigt der Fremdschlüssel, der in dieser Zeile gespeichert ist, auf unsere neue Zeile in Tabelle~\sqlil{g}.%
}%
%
\item<44-> Weil beide Unterausdrücke Teil des selben Kommandos sind, wird die referentielle Integrität -- also die \sqlil{REFERENCES}-Einschränkung -- erst am Ende geprüft, wenn das~\sqlil{;} erreicht wird.%
%
\item<45-> Zu dieser Zeit stimmt es wieder und die referentielle Integrität ist intakt.%
%
\item<46-> Wir haben eine Zeile in Tabelle~\sqlil{g} eingefügt und die passende Zeile in Tabelle~\sqlil{h} referenziert sie auch.%
%
\item<47-> Es ist dann viel einfacher, existierende Zeilen in~\sqlil{g} mit neuen Zeilen in~\sqlil{h} zu verbinden.%
%
\item<48-> Das geht mit einem einfachen \sqlilIdx{UPDATE}~Statement auf Tabelle~\sqlil{h}.
}\end{itemize}%
}%
%
\only<5>{%
\cquotation{PGDG:PD:CT2}{%
A constraint that is not deferrable will be checked immediately after every command. %
Checking of constraints that are deferrable can be postponed until the end of the transaction\dots%
}%
}%
%
\gitLoadAndExecSQL{GH_insert_and_select}{}{conceptualToRelational}{GH_insert_and_select.sql}{relationships}{}{}%
\listingSQLandOutput{-4,6-}{GH_insert_and_select}{}{0.47}{0.01}{0.529}{0.995}
\end{frame}%

\begin{frame}[t]%
\frametitle{\crowsFoot{G}{O1}{H: }{MM}Der Inhalt der Zwei Tabellen}%
%
%
\gitExec{cdtrmTableG}{\databasesCodeRepo}{conceptualToRelational}{../_scripts_/db_table_to_latex_table.sh relationships g gid;fkhid;x}%
\gitExec{cdtrmTableH}{\databasesCodeRepo}{conceptualToRelational}{../_scripts_/db_table_to_latex_table.sh relationships h hid;fkgid;y}%
%
\vfill%
%
\begin{itemize}%
\item Dies sind die Daten in den zwei Tabellen.%
\end{itemize}%
\vfill%
%
\begin{center}%
\strut\hfill\strut%
\centering\input{\gitFile{cdtrmTableG}}%
\strut\hfill\strut%
\centering\input{\gitFile{cdtrmTableH}}%
\strut\hfill\strut%
\end{center}%
%
\hfill%
\end{frame}%
%
%
\begin{frame}[t]%
\frametitle{\crowsFoot{G}{O1}{H: }{MM}Testen der Einschränkungen}%
\parbox{0.435\linewidth}{\small%
\begin{itemize}%
%
\only<-5>{%
\item Können wir eine zweite Zeile in Tabelle~\sqlil{g} auf eine Zeile in Tabelle~\sqlil{h} zeigen lassen, die schon mit einer anderen Zeile in Tabelle~\sqlil{g} verbunden ist?%
}%
%
\only<-6>{%
\item<2-> Unser Modell erlaubt das nicht.%
}%
%
\only<-7>{%
\item<3-> Und es geht auch nicht.%
}%
%
\item<4-> Können wir eine Zeile in Tabelle~\sqlil{h} auf eine andere Zeile in Tabelle~\sqlil{g} verweisen lassen?%
%
\item<5-> Nein, das geht auch nicht.%
%
\item<6-> Können wir eine Zeile in Tabelle~\sqlil{h}, die aktuell mit einer Zeile in Tabelle~\sqlil{g} verbunden ist, auf eine andere Zeile zeigen lassen~(via \sqlil{UPDATE})?%
%
\item<7-> Nein, geht auch nicht.%
%
\item<8-> Unsere Einschränkungen sind ziemlich stark und schützen unsere Daten.%
%
\end{itemize}%
}%
%
\gitLoadAndExecSQL{GH_insert_error_3}{}{conceptualToRelational}{GH_insert_error_3.sql}{relationships}{}{}%
\gitLoadAndExecSQL{GH_insert_error_4}{}{conceptualToRelational}{GH_insert_error_4.sql}{relationships}{}{}%
\gitLoadAndExecSQL{GH_insert_error_5}{}{conceptualToRelational}{GH_insert_error_5.sql}{relationships}{}{}%
\listingSQLandOutput{-3}{GH_insert_error_3}{}{0.47}{0.1}{0.529}{0.87}%
\listingSQLandOutput{4-5}{GH_insert_error_4}{}{0.47}{0.1}{0.529}{0.87}%
\listingSQLandOutput{8-9}{GH_insert_error_5}{}{0.47}{0.1}{0.529}{0.87}%
\end{frame}%
